\documentclass[11pt]{article}
\usepackage[a4paper,margin=1in]{geometry}
\usepackage{amsmath,amssymb,mathtools,bm}
\usepackage{hyperref}
\usepackage{physics}
\usepackage{enumitem}

\title{Checking the Mode Equations in\\ \textit{Cosmological gravitational particle production of massive spin-2 particles}}
\author{Project notes for \texttt{Spin2CGPP}}
\date{\today}

\begin{document}
\maketitle

\section{Scope}
These notes document a reproducible check of the mode equations in Siyang Ling, Edward W. Kolb, Andrew J. Long, and Rachel A. Rosen, \texttt{arXiv:2302.04390}. The target equations are the sector-by-sector equations of motion in section 3 of the paper:
\begin{itemize}[leftmargin=1.5em]
\item minimal matter coupling: tensor, vector, and canonical scalar equations;
\item nonminimal matter coupling: tensor, vector, and scalar equations.
\end{itemize}

The scripts in \texttt{Code/script/} start from the published quadratic sector Lagrangians and reproduce the final mode equations by explicit Euler-Lagrange variation, including the canonical field redefinitions used in the paper.

\section{What the automation does}
The derivation pipeline is organized around a single repeated pattern. If a sector Lagrangian has the form
\begin{equation}
\mathcal{L} = K(\eta)\, q'^{\,2} - M(\eta)\, q^2,
\end{equation}
then after the canonical redefinition
\begin{equation}
q(\eta) = \frac{\chi(\eta)}{\sqrt{2K(\eta)}}
\end{equation}
and after discarding a total derivative, the Lagrangian becomes
\begin{equation}
\mathcal{L} = \frac{1}{2}\chi'^{\,2} - \frac{1}{2}\omega_k^2(\eta)\chi^2,
\qquad
\omega_k^2(\eta) = \frac{4KM + (K')^2 - 2K K''}{4K^2}.
\end{equation}
The Mathematica helper \texttt{CanonicalOmegaSquared} implements exactly this identity.

For the minimal scalar sector, the paper already presents the canonically normalized two-field system
\begin{equation}
\mathcal{L}_S =
\frac12 |\chi_\Pi'|^2 - \frac12 \omega_\Pi^2 |\chi_\Pi|^2
+ \frac12 |\chi_{\mathcal B}'|^2 - \frac12 \omega_{\mathcal B}^2 |\chi_{\mathcal B}|^2
+ \sigma_1 \chi_\Pi^\ast \chi_{\mathcal B}' - \sigma_0 \chi_\Pi^\ast \chi_{\mathcal B}.
\end{equation}
In the Mathematica checks, I model the derivative mixing with an equivalent real-field convention. That convention differs from the paper's complex-mode presentation by a total-derivative reshuffling, so an explicit $\sigma_1'/2$ shift appears in the off-diagonal mass-like term. After accounting for that convention choice, the coupled equations agree with the structure quoted in the paper:
\begin{align}
\chi_\Pi'' + \omega_\Pi^2 \chi_\Pi - \sigma_1 \chi_{\mathcal B}' + \sigma_0 \chi_{\mathcal B} &= 0, \\
\chi_{\mathcal B}'' + \omega_{\mathcal B}^2 \chi_{\mathcal B} + \sigma_1 \chi_\Pi' + \sigma_0 \chi_\Pi &= 0.
\end{align}

\section{Minimal matter coupling}
\subsection{Tensor sector}
The paper starts from
\begin{equation}
\mathcal{L}_T = \frac12 a^2 \left[ D_{ij}' D_{ij}' - \partial_k D_{ij}\partial_k D_{ij} - a^2 m^2 D_{ij}D_{ij}\right].
\end{equation}
After Fourier transforming and replacing $\partial_k \partial_k \to -k^2$, one obtains a single-mode Lagrangian of the form
\begin{equation}
\mathcal{L}_{T,k} = \frac12 a^2 \left[ D'^2 - (k^2 + a^2 m^2) D^2 \right].
\end{equation}
Using $\chi = aD$, the code shows
\begin{equation}
\mathcal{L}_{T,k} = \frac12 \chi'^2 - \frac12 \left(k^2 + a^2 m^2 - \frac{a''}{a}\right)\chi^2
\end{equation}
up to a total derivative. With $H = a'/a^2$, the background identity
\begin{equation}
\frac{a''}{a} = 2a^2 H^2 + aH'
\end{equation}
gives
\begin{equation}
\omega_k^2 = k^2 + a^2 m^2 - 2a^2 H^2 - aH',
\end{equation}
which matches eq.~(3.8) of the paper.

\subsection{Vector sector}
The reduced Fourier-space Lagrangian quoted in the paper is
\begin{equation}
\mathcal{L}_{V,k} = K_C |C'|^2 - M_C |C|^2,
\qquad
K_C = \frac{a^4 k^2 m^2}{k^2 + a^2 m^2},
\qquad
M_C = a^4 k^2 m^2.
\end{equation}
Applying the generic canonical-normalization identity above gives
\begin{equation}
\omega_k^2 =
\frac{4K_C M_C + (K_C')^2 - 2 K_C K_C''}{4K_C^2}.
\end{equation}
The script rewrites this result as
\begin{equation}
\omega_k^2 = k^2 + a^2 m^2 - \frac{f''}{f},
\qquad
f = \frac{a^2}{\sqrt{k^2 + a^2 m^2}},
\end{equation}
which is exactly eq.~(3.13).

\subsection{Scalar sector}
The minimally coupled scalar sector is the technically hardest part of the paper. The current project now reconstructs the scalar reduction directly from the covariant quadratic action. The scripts
\begin{itemize}[leftmargin=1.5em]
\item \texttt{Code/script/scalar\_equation\_rank\_check.wls},
\item \texttt{Code/script/derive\_scalar\_mode\_eoms.wls}, and
\item \texttt{Code/test/test\_scalar\_mode\_chain.wls}
\end{itemize}
implement the scalar SVT ansatz, perform the field redefinition
\begin{equation}
\hat{\varphi}_v = \varphi_v - \frac{a^{-1}\bar\phi'}{\Mpl H} A,
\end{equation}
and then eliminate the constrained variables in the order
\begin{equation}
F_\pm \longrightarrow A_\pm \longrightarrow E_\pm.
\end{equation}

This ordering is important. A naive Hessian analysis of the partially reduced Lagrangian suggests that $A$ remains dynamical, but that conclusion is misleading. After solving the $F_\pm$ constraints, the $E_\pm$ equations solve $A_\pm$ without introducing $A_\pm''$; after substituting those solutions, the $A_\pm$ equations solve $E_\pm$. Only after this full sequential elimination does the expected coupled two-field system in $(B,\hat\varphi_v)$ emerge.

The resulting equations of motion are denoted \texttt{finalB} and \texttt{finalP} in \texttt{derive\_scalar\_mode\_eoms.wls}. The test script verifies three structural facts directly:
\begin{enumerate}[leftmargin=1.5em]
\item both final equations contain $B''$ and $\hat\varphi_v''$;
\item the final $2\times2$ acceleration matrix has rank $2$;
\item the derivation starts from the covariant action rather than the published reduced scalar action.
\end{enumerate}

The current project now includes an explicit Mathematica transcription of the published kinetic-sector coefficients $K_\varphi$, $K_B$, $L_2$, $L_1$, together with the derived combinations $\kappa$ and $K_\mathcal{B}$, in \texttt{Code/script/paper\_scalar\_coeffs.wl}. The regression check \texttt{Code/test/test\_scalar\_paper\_kinetics.wls} verifies that the diagonal acceleration coefficients extracted from the covariant scalar chain match the paper's $K_\varphi$ and $K_B$ formulas after equation normalization, and that the resulting acceleration matrix can be lifted into the paper's kinetic matrix.

What remains unsolved is the full reduced-action matching problem rather than the existence of the scalar equations themselves. The sequentially eliminated equations are not yet in the same reduced-action basis as the paper: after lifting the acceleration matrix into the published $(K_\varphi,L_2;L_2,K_B)$ basis, the first-derivative matrix still shows a nonzero residual. So the remaining task is to reconstruct the correct reduced-action basis and then compare the full coefficient set $(K_\varphi,M_\varphi,K_B,M_B,L_2,L_1,L_0)$ term by term.

\section{Nonminimal matter coupling}
\subsection{Tensor sector}
The tensor-sector action in the nonminimal model is
\begin{equation}
\mathcal{L}_T = \frac12 a^2 \left[D'^2 - \left(k^2 + a^2(m^2-\Lambda+3H^2+2a^{-1}H')\right)D^2\right].
\end{equation}
Again using $\chi = aD$, the canonical form is
\begin{equation}
\mathcal{L}_{T,k} = \frac12 \chi'^2 - \frac12 \omega_k^2 \chi^2
\end{equation}
with
\begin{equation}
\omega_k^2 = k^2 + a^2 m^2 - a^2 \Lambda + a^2 H^2 + aH',
\end{equation}
matching eq.~(3.58).

\subsection{Vector and scalar sectors}
The paper reduces both the vector and scalar sectors to single-field quadratic actions of the same form
\begin{equation}
\mathcal{L} = K(\eta)\,q'^{\,2} - M(\eta)\,q^2.
\end{equation}
The Mathematica scripts apply the same canonical-normalization identity and reproduce the quoted frequency formulas
\begin{align}
\omega_{V,k}^2 &= \frac{4K_C M_C + (K_C')^2 - 2K_C K_C''}{4K_C^2}, \\
\omega_{S,k}^2 &= \frac{4K_B M_B + (K_B')^2 - 2K_B K_B''}{4K_B^2},
\end{align}
which are eqs.~(3.66) and (3.73).

\section{Files}
\begin{itemize}[leftmargin=1.5em]
\item \texttt{Code/script/Spin2CGPP.wl}: derivation helpers and symbolic checks.
\item \texttt{Code/script/derive\_eoms.wls}: human-readable driver that prints each check.
\item \texttt{Code/script/derive\_scalar\_mode\_eoms.wls}: covariant-to-scalar sequential elimination chain for the minimally coupled scalar sector.
\item \texttt{Code/script/check\_ds\_scalar\_limit.wls}: de Sitter-limit diagnostic for the pure spin-2 scalar reduction.
\item \texttt{Code/test/run\_tests.wls}: Mathematica regression tests.
\item \texttt{Code/test/test\_scalar\_mode\_chain.wls}: regression test for the covariant scalar elimination chain.
\item \texttt{Code/test/run\_tests.sh}: shell entry point for the tests.
\end{itemize}

\section{Current limitation}
The largest remaining gap is no longer the covariant-to-scalar derivation itself, nor the diagonal kinetic coefficients. Those pieces are now automated and checked. The remaining gap is the full reduced-action matching of the minimally coupled scalar sector: the current sequentially eliminated equations still need to be converted into the same reduced-action basis used for eq.~(3.18), so that the complete coefficient set $(K_\varphi,M_\varphi,K_B,M_B,L_2,L_1,L_0)$ can be compared term by term. The nonminimal scalar-sector derivation from the covariant action also remains to be automated at the same level.

\end{document}
